% !TeX root = ../papers/nc-p3.tex
\section{頂点指数四で四次中心系の基点になれる単純繊維}
\label{sec:index-four-quartic-basepoints}
\begin{theorem}\label{thm:index-four-quartic-basepoints}
有限四辺の標準正則窓が中心上有限で、頂点可除指数が4とする。非頂点閉中心指標 $\chi$ で
経路長4の全中心元が消える必要十分条件は、次の三つのうち一つが成立することである。
\begin{enumerate}
\item 各頂点四次元の単純一個で、対向する隣接二段の行列積空間がtrace対合で直交する。
\item 二つの各頂点二次元単純で、一方がもう一方の経路長4の作用を $-1$ 倍した次数倍率同型類である。
\item 二つの異なるグラフ台点に属する薄い単純が二個ずつあり、各組の一周倍率が互いに逆符号である。
\end{enumerate}
ここでは一つの $\chi$ の繊維を取って判定する。後二種の指標は、経路長を4で割った中心の次数作用の
位数二の元に固定され、八次中心系の基点ではない。第一種の指標で全中心のアフィン錐は滑らかである。
\end{theorem}
\begin{proof}
一般の各隅の被約traceは同次数の全中心に延び、中心元に制限すれば4倍である。従って一つの始点での
全ての四文字のtraceが零かどうかで四次基点を判定できる。前節の付値格子の議論で、繊維の全ての
単純のtraceが正整数重み付きで現れ、頂点次元の重み付き和は4である。各単純の四文字traceテンソルを
$\tau_S(v_0,v_1,v_2,v_3)=\operatorname{tr}(A_3(v_3)\cdots A_0(v_0))$ とする。

頂点次元三の単純の $\tau_S$ はどの一文字でも階数1以下にはならない。例えば第二位置でそうなら
その四文字空間の超平面の像 $K\subset H_1=A_1(V_1)$ と $H_0=A_0(V_0)$ に対し $KH_0$ は
もう一方の二段積のtrace直交に入る。この直交には階数1がない。実際もう一方の積は各ベクトル上で全射である。
$H_0$ に階数1行列があればその像直線を $K$ が消すことになり、$H_1$ 全体での像次元2以上に反する。
$K$ に階数1があれば同様に $H_0$ の全像がその核に入ってしまう。
$H_i$ の次元は3以上である。二次元以下なら各非零ベクトルで像が二次元以上なので、二次元の
任意の非零行列結合が可逆でなければならず不可能である。従って $K$ には階数2行列 $B$ がある。
$H_0$ のどの階数2行列 $A$ の像平面も $\ker B$ を含まない。含めば $BA$ が階数1になる。
射影的な $\det A=0$ の軌跡は全体または次元2以上の射影空間の超曲面で連結であり、
階数はそこでは常に2である。像平面への射は双対射影平面のアフィン部分（$\ker B$ を避ける平面）に
入るから一定となる。しかし一定平面を消す共変元 $u$ の $uH_0$ は双対評価で次元2以上であり、
$uA=0$ の線型条件がこの全超曲面を含むことはできない。矛盾である。これで $3+1$ は消去できる。
$2+1+1$ も二次元単純の二段積同士のtrace対合の階数4と、二つの純テンソルの階数高々2から消去できる。

二次元単純の $\tau_S$ の射影類はその単純を共通次数倍率まで復元する。一つの文字 $v$ を固定し、
次の文字から残り二文字の双線型式への階数を見る。残り二段の積は全二行二列を張るので
この階数は $\dim H_{i+1}A_i(v)$ であり、$A_i(v)$ の階数が $0,1,2$ の順に $0,2,\dim H_{i+1}$
となる。最後は3または4であるから行列式二次曲面が復元される。階数1文字の次文字側の核は
$\operatorname{im}A_i(v)$ を消す矢の二平面である。違う像直線に同じ核が対応することはない。
両方を消す矢の空間は $\ker A_{i+1}$ で次元1以下だからである。従って滑らかな二次曲面では
像直線による罫線族も復元できる。四位置に定理\ref{thm:finite-two-dimensional-quadrics}を使う。
$2+2$ の消滅は従って全く同じ射影単純の相反する一周倍率二つが各一回出る場合に限り、逆も成立する。

薄い単純だけなら高々四台点に支えられた四文字評価は独立である。さもなければ長さ20の
グラフ上で四積を消す退化積対合が生じる（定理\ref{thm:finite-four-slot-rank}）。各台点で
倍率の重み付き和が零となり、少なくとも二つの異なる倍率型が必要である。一台点のみなら前節の異倍率間の
拡張消滅とブロック連結性に反する。従って第三種だけが残り、逆も明らかである。
第一種では対向積の直交が $\tau_S=0$ そのものである。全折り畳みのPI次数16に達する単純を持つので、
Artin--ProcesiのAzumaya判定でその指標上はAzumayaである。有限大域次元を局所化しエタール局所で行列環にすると
その中心は正則である（または中心有限なhomologically homogeneous環の滑らかなAzumaya軌跡の議論を使う）。

後二種で逆符号の一方の表現に次数作用を施せば、同じ中心指標のもう一方となる。
従ってその作用は $\chi$ を固定する。八次traceが全て消えることはない。二次元の二型なら
四段の積空間は全行列環で、二つの四段の積に対する二型のtraceは同じ符号で加わる。
薄い二組なら八文字で符号が消え、異なる二台点の評価は独立である。
\end{proof}
